% ============================================================
%  Section 4.4 – Tests de conformité
% ============================================================
%
%  Appliquer la méthodologie en 5 étapes (§ 3.2.2 du cours) pour chaque test :
%
%  Étape 1 : Définir H0 et H1 (reprendre § 02_formulation.tex)
%  Étape 2 : Distribution de la statistique sous H0
%             - Proportion : N(0,1) si n≥30, np≥5, n(1−p)≥5
%             - Moyenne    : T(n−1) si n≥30 (variance σ inconnue)
%  Étape 3 : Région critique
%             - Bilatéral  : intervalle [−zα ; +zα] ou [−t_{n−1,α} ; +t_{n−1,α}]
%             - Unilatéral : seuil z2α ou t_{n−1,2α}
%  Étape 4 : Valeur du critère
%             - Proportion : z = (p0 − p) / sqrt(p(1−p)/n)
%             - Moyenne    : t = (x̄ − µ) / (σe/√n)
%  Étape 5 : Conclusion (rejeter ou ne pas rejeter H0)
%
%  Si les conditions ne sont pas réunies : le mentionner explicitement
%  et présenter le test quand même (consigne du professeur).
%
% ============================================================

\subsection{Tests de conformité}

% TODO : appliquer la méthodologie en 5 étapes pour chaque test
